\section{Quantitative multiplicity, optimization and transfer}
\label{sec:quantitative-multiplicity}

Retain a fixed reduced target $t=a/b\in(0,1)$ with squarefree $b$, a fixed admissible exponent $\gamma$, and a fixed forbidden set.  Write $\mathcal E_X$, $L_X$, $\Lambda_X$, $\theta_X$ and $\sigma_X$ for the parameterized objects above.  Let
\[
 M_X=|\mathcal E_X|,
 \qquad
 \mathcal R_X(t)=\left\{A\subseteq\mathcal E_X:
   \sum_{e\in A}\frac1e=t\right\},
\]
and let $P_X(t)$ be the Bernoulli probability of the exact target event.

\subsection{Sharp fixed-target coefficient from the total variance}

By character orthogonality and no-wrap,
\begin{equation}
\label{eq:quantitative-weighted-count}
 P_X(t)=\sum_{A\in\mathcal R_X(t)}
 \theta_X^{|A|}(1-\theta_X)^{M_X-|A|}.
\end{equation}
The basic asymptotics are
\begin{equation}
\label{eq:quantitative-basic-asymptotics}
\begin{aligned}
 \theta_X&\longrightarrow\alpha_{t,\gamma}
 =\frac{2t}{(\log\gamma)^2},\\
 M_X&=(1+o(1))\frac{Z^2}{2\log^2Z},\\
 \log L_X&=(1+o(1))Z=o(M_X).
\end{aligned}
\end{equation}
The variance is the total actual-family variance \cref{eq:total-variance}; its two leading contributions are separated only afterwards by \cref{cor:variance-regimes}.

\begin{corollary}[Quantitative target coefficient]
\label{cor:quantitative-coefficient}
For all sufficiently large $X$,
\begin{equation}
\label{eq:quantitative-coefficient}
 L_XP_X(t)\geq\frac{c_{t,\gamma}}{\sigma_X}
\end{equation}
for a fixed $c_{t,\gamma}>0$.
\end{corollary}

\begin{proof}
The strict six-sector assembly leaves a fixed positive fraction of the major contribution after all minor totals have been subtracted.
\end{proof}

\begin{theorem}[Sharp fixed-target coefficient]
\label{thm:sharp-target-coefficient}
For fixed $t$, $\gamma$ and finite forbidden set,
\begin{equation}
\label{eq:sharp-target-coefficient}
 L_XP_X(t)\sim\frac1{\sqrt{2\pi}\,\sigma_X}.
\end{equation}
Equivalently, with $\alpha=\alpha_{t,\gamma}$,
\begin{equation}
\label{eq:sharp-target-coefficient-total}
 L_XP_X(t)
 \sim
 \left[
 2\pi\alpha(1-\alpha)
 \left(
  \frac1{2X^2\log^2X}
  +\frac{\tau(b)}{Z\log Z}
 \right)
 \right]^{-1/2}.
\end{equation}
The explicit variance regimes give
\begin{equation}
\label{eq:sharp-target-coefficient-regimes}
 L_XP_X(t)\sim
 \begin{cases}
 \displaystyle
 \frac{X\log X}{\sqrt{\pi\alpha(1-\alpha)}},&\gamma>2,\\[8pt]
 \displaystyle
 \frac{\sqrt{Z\log Z}}
 {\sqrt{2\pi\alpha(1-\alpha)\tau(b)}},&1<\gamma\leq2.
 \end{cases}
\end{equation}
\end{theorem}

\begin{proof}
Fix $C\geq1$ while $X\to\infty$.  With $N=\lfloor C/\sigma_X\rfloor$, the uniform complex expansion \cref{eq:major-local-gaussian,eq:aggregate-cubic-remainder} gives
\[
 \sigma_X\sum_{|m|\leq C/\sigma_X}F(m)
 \longrightarrow\int_{-C}^{C}\exp(-2\pi^2u^2)\,du.
\]
By \cref{eq:sector-II-tail}, the normalized total-variance tail is $O(\exp(-cC^2))$.  For fixed $C$, \cref{eq:transition-total,eq:adaptive-total,eq:anchor-tail-total,eq:terminal-error-total} make every remaining sector $o(\sigma_X^{-1})$.  Hence
\[
 \limsup_{X\to\infty}
 \left|\sigma_XL_XP_X(t)-
 \int_{-C}^{C}\exp(-2\pi^2u^2)\,du\right|
 \ll\exp(-cC^2).
\]
Only now let $C\to\infty$.  Since the Gaussian integral is $1/\sqrt{2\pi}$, \cref{eq:sharp-target-coefficient} follows.  Formula \cref{eq:sharp-target-coefficient-total} is the total variance \cref{eq:total-variance}, and the two explicit forms are \cref{cor:variance-regimes}.  No diagonal choice $C=C(X)$ and no moving-target uniformity is used.
\end{proof}

\subsection{Direct-family entropy and support diversity}

Put
\[
 H(u)=-u\log u-(1-u)\log(1-u)
 \qquad(0<u<1),
\]
and let
\[
 K_X\sim\operatorname{Bin}(M_X,\theta_X),
 \qquad
 w_X=\left\lceil2\sqrt{M_X\log L_X}\right\rceil.
\]
The modulus estimate gives
\begin{equation}
\label{eq:wX-scales}
 \frac{w_X}{\sqrt{M_X}}\sim2\sqrt{\log L_X}\longrightarrow\infty,
 \qquad
 \frac{w_X}{M_X}
 =O\!\left(\sqrt{\frac{\log L_X}{M_X}}\right)=o(1).
\end{equation}

\begin{theorem}[Mesoscopic entropy at the direct-family rate]
\label{thm:mesoscopic-entropy}
One has
\begin{equation}
\label{eq:mesoscopic-entropy-lower}
 \#\left\{A\in\mathcal R_X(t):
 \bigl||A|-\theta_XM_X\bigr|\leq w_X\right\}
 \geq\exp\bigl([H(\alpha_{t,\gamma})-o(1)]M_X\bigr).
\end{equation}
The number of all subsets of $\mathcal E_X$ in the same window is at most
\begin{equation}
\label{eq:mesoscopic-entropy-upper}
 \exp\bigl([H(\alpha_{t,\gamma})+o(1)]M_X\bigr).
\end{equation}
Thus the lower bound has the optimal exponential rate within the stated cardinality window.
\end{theorem}

\begin{proof}
Hoeffding's inequality and the definition of $w_X$ give
\begin{equation}
\label{eq:hoeffding-window}
 \Prob(|K_X-\theta_XM_X|>w_X)
 \leq2\exp(-2w_X^2/M_X)\leq2L_X^{-8}.
\end{equation}
The sharp coefficient gives the quantitative comparison that is needed here:
\begin{equation}
\label{eq:hoeffding-target-ratio}
 \frac{2L_X^{-8}}{P_X(t)}
 \sim2\sqrt{2\pi}\,\sigma_XL_X^{-7}.
\end{equation}
By \cref{eq:total-variance},
\[
 \sigma_X=O\!\left(\frac1{X\log X}+\frac1{\sqrt{Z\log Z}}\right)=o(1),
\]
while \cref{prop:family-size-modulus} gives $\log L_X=(1+o(1))Z$ and hence
\[
 L_X^{-7}=\exp(-(7+o(1))Z).
\]
Therefore $\sigma_XL_X^{-7}\to0$, and \cref{eq:hoeffding-target-ratio} proves
\begin{equation}
\label{eq:hoeffding-negligible}
 2L_X^{-8}=o(P_X(t)).
\end{equation}
Thus a $(1-o(1))$ proportion of the exact target mass lies in the window.  The explicit variance regimes also give
\[
 |\log\sigma_X|=O(\log Z)=o(M_X).
\]
Together with $\log L_X=o(M_X)$ and the sharp coefficient, this yields
\begin{equation}
\label{eq:target-probability-subexponential}
 \log P_X(t)=-\log L_X-\log\sigma_X+O(1)=-o(M_X).
\end{equation}

For $k=\theta_XM_X+d$ with $|d|\leq w_X$, compactness of the Bernoulli interval gives
\[
\begin{aligned}
 \log\bigl(\theta_X^k(1-\theta_X)^{M_X-k}\bigr)
 &=-M_XH(\theta_X)+d\log\frac{\theta_X}{1-\theta_X}\\
 &\leq-M_XH(\theta_X)+O(w_X).
\end{aligned}
\]
Divide the target mass in the window by this maximal atom and use \cref{eq:quantitative-basic-asymptotics,eq:wX-scales} to obtain \cref{eq:mesoscopic-entropy-lower}.  For the upper bound, use
\[
 \binom{M_X}{k}\leq\exp\bigl(M_XH(k/M_X)\bigr).
\]
Uniformly in the window, $k/M_X=\theta_X+o(1)$ by \cref{eq:wX-scales}, and the number of possible cardinalities is $2w_X+1=\exp(o(M_X))$.  Continuity of $H$ gives \cref{eq:mesoscopic-entropy-upper}.
\end{proof}

\begin{corollary}[Exact-cardinality extraction]
\label{cor:exact-cardinality-extraction}
There is an integer
\[
 k_X=\theta_XM_X+O(w_X)
\]
such that
\begin{equation}
\label{eq:exact-cardinality-count}
 \#\{A\in\mathcal R_X(t):|A|=k_X\}
 \geq\exp\bigl([H(\alpha_{t,\gamma})-o(1)]M_X\bigr).
\end{equation}
\end{corollary}

\begin{proof}
Pigeonhole among the $2w_X+1=\exp(o(M_X))$ cardinalities in the window.
\end{proof}

\begin{corollary}[Macroscopic diversity at one cardinality]
\label{cor:hamming-separated-family}
There are constants $\rho,\eta>0$ and, for all sufficiently large $X$, a family $\mathcal F_X\subseteq\mathcal R_X(t)$ such that
\[
 |\mathcal F_X|\geq\exp(\eta M_X),
 \qquad
 |A|=k_X\quad(A\in\mathcal F_X),
\]
and
\[
 |A\mathbin\triangle B|\geq\rho M_X
 \qquad(A\neq B\in\mathcal F_X).
\]
\end{corollary}

\begin{proof}
Choose $\rho>0$ so small that $H(\rho)<H(\alpha_{t,\gamma})$.  A Hamming ball of radius $\rho M_X$ has at most $\exp([H(\rho)+o(1)]M_X)$ members.  Greedy packing inside the exact-cardinality family proves the assertion.
\end{proof}

\begin{corollary}[Balanced signed reciprocal relations]
\label{cor:balanced-signed-relations}
There are constants $c,\eta'>0$ such that, for all sufficiently large $X$, there are at least $\exp(\eta'M_X)$ distinct ordered pairs $(\mathcal P,\mathcal N)$ of disjoint subsets of $\mathcal E_X$ satisfying
\[
 \sum_{e\in\mathcal P}\frac1e
 =\sum_{e\in\mathcal N}\frac1e,
 \qquad
 |\mathcal P|=|\mathcal N|\geq cM_X.
\]
\end{corollary}

\begin{proof}
Fix $A_0\in\mathcal F_X$.  For every other $B\in\mathcal F_X$, take $\mathcal P=A_0\setminus B$ and $\mathcal N=B\setminus A_0$.  Equal cardinality and equal reciprocal target give the two equalities, while Hamming separation gives the linear support size.  The map from $B$ is injective.
\end{proof}

